\documentclass[aspectratio=169]{beamer}

% Packages for Ukrainian language support
\usepackage[T2A]{fontenc}
\usepackage[utf8]{inputenc}
\usepackage[ukrainian]{babel}

% Theme and colors
\usetheme{Madrid}
\usecolortheme{default}

% Custom colors
\definecolor{chatgptgreen}{RGB}{16, 163, 127}
\setbeamercolor{structure}{fg=chatgptgreen}
\setbeamercolor{title}{fg=white, bg=chatgptgreen}
\setbeamercolor{frametitle}{fg=white, bg=chatgptgreen}

% Title page information
\title{Знайомтесь: ChatGPT}
\subtitle{Ваш новий цифровий помічник}
\author{}
\date{}

\begin{document}

\begin{frame}
    \titlepage
\end{frame}

\begin{frame}{Слайд 1: Що таке ChatGPT?}
    \begin{itemize}
        \item \textbf{Концепція:} Уявіть, що у вас є дуже розумний, начитаний помічник, який не спить 24/7 і готовий допомогти вам майже з чим завгодно. Це комп'ютерна програма, з якою ви можете розмовляти, як з людиною.
        \item \textbf{Приклад з реального життя:} Замість того, щоб шукати в Google і читати 10 різних веб-сайтів, щоб знайти відповідь, ви просто ставите запитання ChatGPT, і він одразу пише для вас точну відповідь.
        \item \textbf{Спробуйте цей запит:} \textit{"Я планую 3-денну поїздку до Відня. Я люблю класичну музику та відвідувати старі кафе. Чи можеш ти скласти для мене маршрут на кожен день?"}
    \end{itemize}
\end{frame}

\begin{frame}{Слайд 2: Зберігаємо свіжість (Чому важливі нові чати)}
    \begin{itemize}
        \item \textbf{Концепція:} Як і люди, якщо ви говорите про занадто багато різних речей в одній довгій розмові, ChatGPT може "виснажитися" або заплутатися. Він може переплутати ваші попередні теми.
        \item \textbf{Правило:} Завжди починайте \textbf{Новий чат} для нової теми. Якщо ви говорили про рецепт, а тепер хочете поговорити про рок-гурт, натисніть кнопку "Новий чат"!
        \item \textbf{Приклад з реального життя:} Якщо у вас є чат, де ви просите допомоги в написанні офіційного електронного листа своєму босу, не просіть раптом рецепт піци в тому ж чаті. Почніть новий!
    \end{itemize}
\end{frame}

\begin{frame}{Слайд 3: Ваш особистий перекладач}
    \begin{itemize}
        \item \textbf{Концепція:} Ви можете перетворити ChatGPT на спеціалізованого, блискавичного перекладача.
        \item \textbf{Як це зробити:} Почніть новий чат і дайте йому одну єдину інструкцію на початку.
        \item \textbf{Запит для копіювання:} \textit{"Відтепер усе, що я пишу українською, перекладай англійською. А все, що я пишу англійською, перекладай українською. Не додавай жодного іншого тексту, просто давай мені переклад."}
        \item \textbf{Магія:} Вам більше ніколи не доведеться просити його перекласти в цьому конкретному чаті. Просто відкрийте цей чат, введіть свій текст, і він автоматично перекладатиме його щоразу!
    \end{itemize}
\end{frame}

\begin{frame}{Слайд 4: Пошук в Інтернеті та пошук посилань}
    \begin{itemize}
        \item \textbf{Концепція:} ChatGPT — це не просто мозок; він також може шукати для вас інформацію в живому Інтернеті.
        \item \textbf{Приклад з реального життя:} Ви хочете купити квитки на концерт, але не знаєте, де шукати або який офіційний веб-сайт.
        \item \textbf{Спробуйте цей запит:} \textit{"Знайди мені офіційний веб-сайт для купівлі квитків на майбутній концерт Coldplay у Варшаві та дай мені пряме посилання."}
    \end{itemize}
\end{frame}

\begin{frame}{Слайд 5: Бачити світ через фотографії}
    \begin{itemize}
        \item \textbf{Концепція:} У ChatGPT є "очі"! Ви можете завантажувати фотографії зі свого телефону та ставити запитання про них.
        \item \textbf{Приклад з реального життя:} Ви гуляєте в парку і бачите красиву, незвичайну квітку, або ви в магазині і бачите гаджет, який не розумієте.
        \item \textbf{Як це зробити:} Зробіть фотографію, завантажте її в ChatGPT і запитайте:
        \item \textbf{Спробуйте цей запит:} \textit{"Що це за квітка на фотографії? Як часто її потрібно поливати, якщо я триматиму її у своїй квартирі?"}
    \end{itemize}
\end{frame}

\begin{frame}{Слайд 6: Уважний слухач (Психологічна підтримка)}
    \begin{itemize}
        \item \textbf{Концепція:} Іноді ми всі відчуваємо смуток, стрес або просто потребуємо з кимось поговорити.
        \item \textbf{Приклад з реального життя:} У вас був дуже напружений день на роботі, і ви просто хочете виговоритися, або ви відчуваєте тривогу через майбутню подію.
        \item \textbf{Спробуйте цей запит:} \textit{"У мене сьогодні був справді жахливий день, і я відчуваю сильний стрес. Чи можеш ти просто вислухати мене і запропонувати кілька втішних слів?"}
    \end{itemize}
\end{frame}

\begin{frame}{Слайд 7: Помічник на кожен день (Кулінарія, здоров'я та попередження!)}
    \begin{itemize}
        \item \textbf{Концепція:} Це фантастичний помічник у повсякденному житті. Ви можете сказати йому, які інгредієнти є у вашому холодильнику, і він дасть вам покроковий рецепт.
        \item \textbf{Приклад з реального життя:} Ви відкриваєте холодильник і бачите лише курку, рис і трохи моркви.
        \item \textbf{Спробуйте цей запит:} \textit{"У мене на кухні є лише куряча грудка, білий рис і морква. Яку просту і смачну вечерю я можу приготувати лише з цих інгредієнтів?"}
        \item \textbf{Попередження:} Хоча він може давати загальні поради щодо здоров'я або самопочуття, \textbf{ніколи не довіряйте ШІ сліпо в серйозних питаннях}. Це не справжній лікар. Завжди перевіряйте важливу інформацію про здоров'я, медицину або безпеку у фахівця-людини.
    \end{itemize}
\end{frame}

\begin{frame}{Слайд 8: Найкращий вчитель (Математика, наука та технології)}
    \begin{itemize}
        \item \textbf{Концепція:} ChatGPT блискуче пояснює складні речі дуже простою мовою.
        \item \textbf{Приклад з реального життя:} Ви чуєте, як люди говорять про "Блокчейн" або "Квантову фізику" в новинах, і ви не маєте уявлення, що це означає.
        \item \textbf{Спробуйте цей запит:} \textit{"Поясни мені, що таке Чорна діра, але поясни це простими словами, ніби я 10-річна дитина."}
    \end{itemize}
\end{frame}

\begin{frame}{Слайд 9: Давайте спробуємо!}
    \begin{itemize}
        \item \textbf{Концепція:} Найкращий спосіб навчитися — це робити. Давайте відкриємо додаток і спробуємо один із запитів прямо зараз!
    \end{itemize}
\end{frame}

\end{document}